\chapter{Modelos de Escolha Múltipla}
\label{cap:multinomial}
% ── índice ──────────────────────────────────────────────────────
\index{logit multinomial|textbf}
\index{mlogit@\texttt{mlogit()}|textbf}
\index{logit ordenado|textbf}
\index{ologit@\texttt{ologit()}|textbf}
\index{probit ordenado|textbf}
\index{oprobit@\texttt{oprobit()}|textbf}
\index{IIA, axioma de|textbf}
\index{independência de alternativas irrelevantes|see{IIA, axioma de}}
\index{categoria de referência}

% ─────────────────────────────────────────────────────────────────────────────
\section{Resposta ordinal e nominal}

Quando o resultado assume mais de duas categorias, duas estruturas são
possíveis:

\begin{itemize}
  \item \textbf{Resposta ordinal}: categorias com ordem natural (muito insatisfeito,
  neutro, muito satisfeito; nível educacional). Modelos ordenados (logit/probit)
  preservam essa estrutura.
  \item \textbf{Resposta nominal}: categorias sem ordem (escolha de modal de
  transporte; tipo de emprego). Logit multinomial sem restrição de ordem.
\end{itemize}

\subsection*{Logit e Probit Ordenados}

O modelo supõe uma variável latente $y^* = x'\beta + \varepsilon$ e cortes
$c_1 < c_2 < \cdots < c_{J-1}$:

\[
  y = j \quad\Longleftrightarrow\quad c_{j-1} < y^* \leq c_j
\]

O logit ordenado usa $\varepsilon \sim \mathrm{Logistic}(0,1)$; o probit
ordena usa $\varepsilon \sim N(0,1)$. Ambos estimam $\beta$ e os cortes por MLE.

\subsection*{Logit Multinomial}

Para $J$ alternativas, com alternativa $j=J$ como base:

\[
  P(y_i = j \mid x_i) = \frac{\exp(\alpha_j + \beta_j'\,x_i)}{1 + \sum_{k=1}^{J-1} \exp(\alpha_k + \beta_k'\,x_i)}
  \quad j = 1,\ldots,J-1
\]

Os coeficientes $(\alpha_j, \beta_j)$ medem o efeito de $x$ sobre as chances
relativas de $j$ versus a alternativa base.

% ─────────────────────────────────────────────────────────────────────────────
\section{Verificação por DGP}

\begin{lstlisting}[caption={DGP para modelos ordenados e multinomial}]
seed(42)
generate df x     = rnormal()
generate df e     = (rnormal() - 0.5) * 1.7321
generate df ystar = 1.5 * x + e
// Categorias 1, 2, 3 com cortes em 0.0 e 1.0
generate df y_ord = 1 + (ystar >= 0.0) + (ystar >= 1.0)

// Multinomial via inverse-CDF
generate df denom  = 1.0 + exp(-1.0 + 2.0*x) + exp(0.5)
generate df p1     = 1.0 / denom
generate df p12    = (1.0 + exp(-1.0 + 2.0*x)) / denom
generate df um     = rnormal()
generate df y_mlog = 1 + (um >= p1) + (um >= p12)
\end{lstlisting}

\subsection*{Ordered Logit}

\begin{lstlisting}
let m_olog = ologit(y_ord ~ x, df)
display m_olog
\end{lstlisting}

\begin{lstlisting}[language={}, basicstyle=\ttfamily\small,
  frame=none, numbers=none, backgroundcolor=\color{codbg},
  xleftmargin=0pt, xrightmargin=0pt, breaklines=false]
=========================== Ordered Logit Results ============================
N=500  Log-L=-396.4473  Pseudo R²=0.1878
-------------------------------- Coefficients --------------------------------
x        4.8137  (0.4633)  z=10.39  ***
Thresholds:  cut1=-0.1472  cut2=3.2191
\end{lstlisting}

\subsection*{Ordered Probit}

\begin{lstlisting}
let m_opro = oprobit(y_ord ~ x, df)
display m_opro
\end{lstlisting}

\begin{lstlisting}[language={}, basicstyle=\ttfamily\small,
  frame=none, numbers=none, backgroundcolor=\color{codbg},
  xleftmargin=0pt, xrightmargin=0pt, breaklines=false]
=========================== Ordered Probit Results ===========================
N=500  Log-L=-392.8369  Pseudo R²=0.1952
-------------------------------- Coefficients --------------------------------
x        2.8683  (0.2603)  z=11.02  ***
Thresholds:  cut1=-0.0768  cut2=1.9076
\end{lstlisting}

A razão entre coeficientes logit e probit é $4{,}81/2{,}87 \approx 1{,}68$,
próximo de $\pi/\sqrt{3} \approx 1{,}81$ — a relação teórica entre as duas
escalas.

\subsection*{Logit Multinomial}

\begin{lstlisting}
let m_mlog = mlogit(y_mlog ~ x, df)
display m_mlog
\end{lstlisting}

\begin{lstlisting}[language={}, basicstyle=\ttfamily\small,
  frame=none, numbers=none, backgroundcolor=\color{codbg},
  xleftmargin=0pt, xrightmargin=0pt, breaklines=false]
==================== Multinomial Logit Regression Results ====================
N=500  Log-L=-513.7171  Pseudo R²=0.0448  AIC=1035.43
------------------------------ y=1 vs base y=3 -------------------------------
const   -0.0809  (0.2024)
x       -0.9006  (0.4069)  *
------------------------------ y=2 vs base y=3 -------------------------------
const   -1.5279  (0.2542)  ***
x        2.0825  (0.4119)  ***
\end{lstlisting}

Para $y=2$ versus base $y=3$: $\hat\beta_x = 2{,}08$ (DGP verdadeiro: $\beta_{x,2} - \beta_{x,3} = 2{,}0$). O sinal negativo de $x$ em $y=1$ vs $y=3$ reflete que $x$ alto aumenta as chances de $y=2$ (que tem $e^{0{,}5}$ de probabilidade base) relativo a $y=1$.

% ─────────────────────────────────────────────────────────────────────────────
\section{Sintaxe}

\begin{lstlisting}
ologit(y ~ x1 + x2, df)
oprobit(y ~ x1 + x2, df)
mlogit(y ~ x1 + x2, df)
mlogit(y ~ x1 + x2, df, base=1)    // alternativa base explícita
\end{lstlisting}

\begin{nota}
  O logit multinomial impõe a propriedade de \textbf{independência de
  alternativas irrelevantes} (IIA): a razão de probabilidades entre duas
  alternativas não depende das demais. Para dados de painel de escolha com
  alternativas correlacionadas, use \lstinline{clogit} (cap.~\ref{cap:painel2}).
\end{nota}

\section{Validação empírica}

O dataset \texttt{beauty} do Wooldridge é usado no caso de validação
\texttt{ologit_beauty} em \texttt{validation/cases/}. O caso estima um logit
ordenado de \hayashi{} e o compara com \texttt{MASS::polr} em R e com
\texttt{statsmodels.miscmodels.ordinal_model} em Python. Execute
\lstinline{hay validate} para ver o estado atual.
